\documentclass[12pt, a4paper]{article}

% Paquetes esenciales para español, matemáticas y gráficos
\usepackage[utf8]{inputenc}
\usepackage[spanish, es-tabla]{babel}
\usepackage{amsmath}
\usepackage{graphicx}
\usepackage{hyperref}
\usepackage{geometry}
\geometry{a4paper, margin=2.5cm}

% Información del documento
\title{\textbf{Memoria: Arquitectura y Modelado}}
\author{David Cendejas Rodríguez y Ángel Jiménez Izquierdo \\ \textit{Borrador para reunión con experto}}
\date{\today}

\begin{document}

\maketitle

\begin{abstract}
Este documento detalla la justificación teórica, el tratamiento de datos y la estructura interna utilizada en el desarrollo de la red neuronal predictiva para la inferencia de parámetros de síntesis FM.
\end{abstract}

\tableofcontents
\newpage

\section{Arquitectura y Modelado de la Red Neuronal}

\subsection{Exploración de Paradigmas: Clasificación vs. Regresión}
\subsubsection{Enfoque de Clasificación}
La clasificación es el paso previo natural a la búsqueda de un modelo de inferencia de parámetros.
De esta manera, se puede comprobar de forma práctica que la red neuronal procesa el audio adecuadamente
y es capaz de conservar su información acústica fundamental.

El \textbf{Prototipo 1} es una prueba de concepto para la primera aproximación al aprendizaje automático aplicado al audio.
En él se utiliza un dataset reducido, construido con la librería NumPy, de alrededor de 50 archivos \textit{.wav} de un segundo de duración.
Estos audios se generan sintéticamente
con frecuencias aleatorias y formas de onda básicas (\textit{sine}, \textit{square}, \textit{sawtooth}, \textit{triangle} y \textit{noise}).

En la fase de preprocesamiento y entrenamiento, desarrollada utilizando TensorFlow, se opta por
no aportar a la red formas de audio crudo. En su lugar, se extraen las características calculando
los Coeficientes Cepstrales en las Frecuencias de Mel (MFCC) de cada muestra. El modelo tiene el objetivo de clasificar 
y asociar la matriz de coeficientes con la etiqueta correspondiente a su tipo de onda. A pesar de esto, 
al plantear una arquitectura básica no convolucional, la red neuronal carece de los mecanismos para
realizar una interpretación correcta de los patrones bidimensionales de los MFCC. Como consecuencia, el 
modelo arroja una precisión muy baja. 

El \textbf{Prototipo 2} se basa en el \textbf{Prototipo 1}. Sigue atendiendo al objetivo de la clasificación, pero tiene 
características que difieren del anterior. Esta vez, los audios se convierten a espectrogramas \textit{.png}
mediante el uso de la librería de \textit{librosa}. Esto conlleva un menor coste computacional y de almacenamiento.
Al tratarse realmente de imágenes, se pueden utilizar modelos preentrenados de visión, encargados de clasificar imágenes.

Para el entrenamiento, se hace uso de \textbf{ResNet34 preentrenada} vía \textit{fastAI}. Para su adaptación, se utilizan \textbf{DataBlocks},
una herramienta de \textit{fastAI} mediante la cual se especifican que las entradas son las imágenes (espectrogramas) y las salidas
son las categorías (tipo de onda). Asimismo, se etiqueta obteniendo la categoría a partir del nombre archivo y se divide
destinando un 80\% de las imágenes para entrenar y un 20\% para validación.

Como resultado de utilizar una red convolucional y la arquitectura externa \textbf{ResNet34}, el modelo obtuvo una buena precisión, especialmente 
identificando ondas puras dentro del rango entrenado. 

\subsubsection{Enfoque de Regresión}
Al plantear un salto hacia la síntesis paramétrica continua, fue necesario tomar una decisión
de diseño fundamental: la elección del motor de síntesis. Tras la evaluación de distintas alternativas, 
se optó por centrar el proyecto en la síntesis por Modulación de Frecuencia (FM). Esta técnica presenta
ventajas claras para el aprendizaje automático. A diferencia de la síntesis aditiva, que requiere decenas
de osciladores simultáneos, la síntesis FM es muy eficiente, consiguiendo generar timbres complejos
a partir de un número reducido de parámetros. De esta forma, se puede comenzar utilizando, por ejemplo, 
la frecuencia portadora, el ratio y el índice de modulación, y luego ampliar el número de variables para
conseguir sonidos más complejos. Estas características acotan la dimensionalidad del problema predictivo, 
haciéndolo viable para la red neuronal.   

Si bien con la clasificación queda validada la extracción de características acústicas, una vez tomada 
esta última decisión, deja de ser útil. Para la inferencia de parámetros de la síntesis FM, resulta 
obligatoria la evolución hacia una arquitectura de regresión.

Para el \textbf{Prototipo 3} se genera un dataset mediano de aproximadamente 15000 muestras, construido 
mediante un barrido de parámetros en un sintetizador FM de \textbf{pyo}. Para la reducción de la carga computacional
y la optimización del espacio, se sustituye \textbf{librosa} por la librería de \textit{torchaudio}. Esta permite
transformar los archivos de audio (\textit{.wav}) en tensores de espectrogramas (\textit{.pt}),
asociados mediante el uso de un (\textit{.csv}) a sus valores numéricos. El 80\% del conjunto de datos se destina 
a entrenamiento y el otro 20\% a validación.

El entrenamiento se implementa mediante \textbf{PyTorch}. Se diseña una arquitectura convolucional 
propia (SmallCNNRegressor). La red neuronal se compone de dos bloques principales:
uno de extracción de características y una cabeza de regresión. 

El primer bloque recibe los espectrogramas en forma de tensores de un único canal (lo que equivaldría a 
imágenes en escala de grises) y los procesa en otros tres bloques convolucionales. Cada uno de estos aplica
kernels de tamaño 3x3 (o núcleos convolucionales) que se deslizan sobre la matriz del espectrograma
para extraer características espaciales, detectando así patrones visuales en la representación espectral del sonido.
A esto le sigue una función de activación no lineal \textbf{ReLU} y una capa de reducción espacial (\textbf{Max Pooling}). 
El último bloque convolucional termina en una capa de agrupamiento (\textbf{Adaptive Average Pooling}), que comprime los
mapas de características a una matriz 4x4.

El siguiente paso es la cabeza de regresión, en donde la información espacio-temporal del espectrograma
(extraída en el bloque anterior), se pasa a un vector de una única dimensión y atraviesa dos capas
completamente conectadas (\textbf{Perceptrón Multicapa}). La primera, comprime la información a 128 neuronas,
y la segunda, saca el vector por una dimensión de salida de tamaño 3, que se corresponde con los tres parámetros
del sintetizador: frecuencia portadora, ratio y el índice de modulación.

El modelo optimiza el proceso de los espectrogramas mediante el algoritmo Adam, utilizando el Error Cuadrático
Medio (MSELoss) como función de pérdida. Y, en cuanto al ciclo de entrenamiento, este se configuró con 
una duración de 5 épocas (\textit{epochs}). En este prototipo se decidió no hacer una búsqueda extensa
de hiperparámetros al tratarse de un prototipo intermedio, con un fin de experimentación inicial 
con la regresión.

A pesar de esto, se puede observar una precisión muy limitada en los resultados experimentales.
Se consigue mantener un MSE bajo en el \textit{Ratio} e \textit{Index}, pero el error en el
\textit{Carrier} resulta desproporcionado, debido a operar en la escala de miles de hercios. 
De este prototipo se puede sacar la conclusión de que el minimizar el error matemático sobre los
parámetros no garantiza la similitud acústica del sonido generado. 


\subsection{El Problema de la Inyectividad y la Necesidad de una Arquitectura Compleja}

\subsection{Pipeline de Datos y Preprocesamiento del Dataset}


\subsection{Diseño de la Arquitectura Propuesta (CNN de Doble Cabeza)}
\subsubsection{Codificador (Encoder)}
\subsubsection{Rama de Inferencia Numérica (Parameter Head)}
\subsubsection{Rama de Reconstrucción (Decoder Head)}

\subsection{Formulación de la Función de Pérdida Híbrida}
\subsubsection{Pérdida en el Dominio de Parámetros (SmoothL1)}
\subsubsection{Pérdida en el Dominio Espectral (L1 + Convergencia Espectral)}

\subsection{Justificación del Stack Tecnológico y Librerías Utilizadas}

\end{document}